\section{部署与复现保障}

\subsection{部署准备 checklist}

部署前必须检视 prompt stability、token budget、fallback 方案与监控 hook。建议用 multi-column table 显示各项指标当前状态，并把每次检查的 artifact 写入 `deploy/logs/`。

\begin{table}[H]
\centering
\begin{tabular}{p{4cm}p{5cm}p{3cm}}
\toprule
检查项 & 验证方式 & 状态 \\
\midrule
Prompt stability & 多 prompt + CoT + self-consistency & 监听中 \\
Token budget & latency + FLOPs shell & 尺寸对齐 \\
Fallback & offline fallback script + human review & 预置 \\
Monitoring hook & loss/latency/attention drift & hook enabled \\
\bottomrule
\end{tabular}
\caption{部署前 checklist。}
\end{table}

\begin{importantbox}{准备 checklist 的使用}
每次部署前用 standardized template 走一遍：prompt log → evaluation dashboard → fallback script。如果某项未通过必须记录在 incident grade-book 并延迟上线。
\end{importantbox}

\subsection{可观测性与复现 tooling stack}

构建 tooling stack 时需包含 instrumentation、LLM plan 追踪、工具调用 trace（doc search、python executor 轮询）和 attention/scoreboard snapshot。

\begin{figure}[H]
\centering
\includegraphics[width=0.85\textwidth]{slides-images/slide-37.jpg}
\caption{Deployment + observability stack（slide-37）。}
\end{figure}

\begin{table}[H]
\centering
\begin{tabular}{p{3.5cm}p{5cm}}
\toprule
层级 & 工具/实践 \\
\midrule
Prompt & Prompt log + versioned dashboard \\
Inference & Monitoring hook + latency/failure metrics \\
Tooling & API call trace + sandbox constraints \\
Observability & Attention/Config snapshot + drift alert \\
\bottomrule
\end{tabular}
\caption{部署部署可观测性 stack。}
\end{table}

\begin{warningbox}{缺失 observability 的风险}
\begin{itemize}
    \item 没有 attention trace，无法确认 emergent jump 是否真；
    \item 没有 tool call log，低层 skill 无法追溯；
    \item 缺乏 drift alert，将延迟 incident trigger。
\end{itemize}
\end{warningbox}

\subsection{本章小结}

部署/复现部分是实践的护栏：从 checklist 到 observability stack，再到 incident grade-book，aim 是让 high-level insights 可落地。

